\begin{longtable}{p{0.15\textwidth} p{0.12\textwidth} p{0.12\textwidth} p{0.12\textwidth} p{0.35\textwidth}}
\caption{The PED Comparative Lexicon: 50 Rigorous Etymologies} \\
\toprule
\textbf{PED Root} & \textbf{Vedic} & \textbf{Old Tamil} & \textbf{Old Tibetan} & \textbf{Scholarly Analysis} \\
\midrule
\endfirsthead
\toprule
\textbf{PED Root} & \textbf{Vedic} & \textbf{Old Tamil} & \textbf{Old Tibetan} & \textbf{Scholarly Analysis} \\
\midrule
\endhead
\bottomrule
\endfoot
\textbf{*NGA} & aham & yan & nga & \small The PED root *nga demonstrates the 'Radial Nasal Shift'. The velar nasal [ŋ] is preserved in Old Tibetan 'nga', but undergoes fronting to the palatal nasal [ɲ] in Old Tamil 'yāṉ' (via *ñāṉ) and correlates with the nasalized ending in Vedic 'aham'. This confirms a 1st person singular nasal onset common to the macro-family. \\ 
\addlinespace[10pt]
\textbf{*TI} & tvam & ni & khyod & \small Root *ti follows the 'Dental-Palatal Invariant'. The dental onset is preserved in Tamil 'nī' (with nasalization) and Vedic 'tvam' (t-), while Tibetan 'khyod' shows a complex cluster where the dental coda -d represents the original invariant. The shift T > n in Dravidian 2nd person is a known diagnostic trait. \\ 
\addlinespace[10pt]
\textbf{*NAM} & vayam & nam & nged & \small PED *nam exhibits 'Radial Nasal Shift'. Tamil 'nām' preserves the original dental-nasal onset and bilabial coda. Vedic 'vayam' shows a shifted onset but retains the nasal coda. Tibetan 'nged' demonstrates the velarization [n > ŋ] characteristic of the Eastern node of the radiation. \\ 
\addlinespace[10pt]
\textbf{*TI-} & idam & itu & di & \small The deictic *ti- follows the 'Dental-Palatal Invariant'. The dental [t] is preserved in Tamil 'itu', while Tibetan 'di' shows voicing. Vedic 'idam' utilizes the same dental base [d] with an added nasal suffix. This proximal deictic is extremely stable across the three families. \\ 
\addlinespace[10pt]
\textbf{*TA-} & tat & atu & de & \small Distal deictic *ta- shows the 'Dental-Palatal Invariant'. Vedic 'tat' and Tamil 'atu' preserve the dental [t], while Tibetan 'de' reflects the voiced variant. This alignment suggests a fundamental tripartite deictic system in PED based on dental onsets. \\ 
\addlinespace[10pt]
\textbf{*KA} & ka & yar & su & \small PED *ka follows 'Velar Radiation'. Vedic 'ka' preserves the velar [k]. Tamil 'yār' reflects a palatalization of the velar initial (*k > c > y), a common Dravidian shift. Tibetan 'su' represents a further weakening to a sibilant, marking the interrogative boundary. \\ 
\addlinespace[10pt]
\textbf{*MA} & na & alla & ma & \small Negative particle *ma. Tibetan 'ma' is the primary reflex. Tamil 'alla' (not so) is a secondary form, but the nasal negation survives in prohibitive contexts. Vedic 'na' shows a nasal shift (m > n), illustrating the 'Radial Nasal Shift' across functional morphemes. \\ 
\addlinespace[10pt]
\textbf{*PA-N} & bahu & pala & mang & \small Root *pa-n demonstrates 'Labial Radiation' (*P > p/b). Tamil 'pala' preserves [p], Vedic 'bahu' shows [b] (with aspiration), and Tibetan 'mang' shows the nasalized coda [ŋ], linking the concept of 'fullness' to the 'Radial Nasal Shift'. \\ 
\addlinespace[10pt]
\textbf{*KA-N} & eka & onru & gcig & \small PED *ka-n for 'one'. Vedic 'eka' preserves the velar [k]. Tamil 'onru' (from *on-tu) preserves the nasal coda. Tibetan 'gcig' (from *k-tyig?) preserves the velar initial. The skeletal alignment [K-N] is robust across all three nodes. \\ 
\addlinespace[10pt]
\textbf{*TI-N} & dva & irantu & gnyis & \small PED *ti-n. 'Dental-Palatal Invariant' [t/d] is clear in Vedic 'dva' and Tamil 'iraṇṭu' (where -ṭ- is a retroflexed dental). Tibetan 'gnyis' shows a palatalized nasal [ny], fitting the 'Radial Nasal Shift' in the coda position. \\ 
\addlinespace[10pt]
\textbf{*TAM} & tri & munru & gsum & \small PED *tam. Vedic 'tri' preserves the dental onset [t]. Tamil 'mūṉṟu' (from *mu-nt-?) and Tibetan 'gsum' show the bilabial nasal [m], either as onset or coda, illustrating the nasal radiation in numerical systems. \\ 
\addlinespace[10pt]
\textbf{*NGA} & panca & aintu & lnga & \small PED *nga. Tibetan 'lnga' is the perfect reflex of the velar nasal. Tamil 'aintu' reflects the nasalized core (from *cay-nt-?). Vedic 'panca' is an IE innovation but retains the nasal [n], possibly influenced by a substrate PED layer. \\ 
\addlinespace[10pt]
\textbf{*PA-R} & mahant & peru & che & \small PED *pa-r 'Great'. Tamil 'peru' preserves the labial onset and liquid coda. Tibetan 'che' (from *khye?) shows velarization. Vedic 'mahant' is distinct but 'pṛthu' (broad) provides a closer cognate for the [P-R] skeleton. \\ 
\addlinespace[10pt]
\textbf{*RA-NG} & dirgha & netu & ring & \small PED *ra-ng. Tibetan 'ring' preserves the [R-NG] skeleton. Vedic 'dīrgha' preserves the liquid and velar (r-gh). Tamil 'neṭu' shows a shift but correlates with the 'Dental-Palatal Invariant' in its suffixal morphology. \\ 
\addlinespace[10pt]
\textbf{*PA-R} & guru & paru & lci & \small PED *pa-r. Tamil 'paru' (thick/heavy) aligns with Vedic 'guru' (where g < *p is rare, but here likely *gwr-). Tibetan 'lci' reflects a complex cluster [l-kyi], preserving the velar radiation of the root. \\ 
\addlinespace[10pt]
\textbf{*KI-R} & alpa & ciru & chung & \small PED *ki-r. Tamil 'ciṟu' shows the palatalization of the velar initial (*k > c). Tibetan 'chung' shows the velar nasalization. Vedic 'alpa' is likely a distinct root, but 'kṣudra' provides a better velar match. \\ 
\addlinespace[10pt]
\textbf{*MA} & matr & tay & ma & \small Nursery root *ma. Universally preserved in Tibetan 'ma' and Vedic 'mātṛ'. Tamil 'tāy' is a distinct honorific but 'amma' remains the primary colloquial and nursery form, confirming the stability of the labial nasal. \\ 
\addlinespace[10pt]
\textbf{*PA} & pitr & tantai & pha & \small Nursery root *pa. Tibetan 'pha' and Vedic 'pitṛ' preserve the labial onset. Tamil 'tantai' is likely a dental-based honorific (*ta-), but 'appa' remains as the labial correlate in the Dravidian substrate. \\ 
\addlinespace[10pt]
\textbf{*NYA} & matsya & min & nya & \small PED *nya. Tibetan 'nya' preserves the palatal nasal. Tamil 'mīṉ' shows the shift to dental nasal [n] while retaining the nasal onset [m]. Vedic 'matsya' preserves the initial labial nasal, following the 'Radial Nasal Shift'. \\ 
\addlinespace[10pt]
\textbf{*PU-L} & vi & pul & bya & \small PED *pu-l. Tamil 'puḷ' is the perfect reflex. Tibetan 'bya' shows the labial-glide shift (*p-y). Vedic 'vi' (from *pwi?) preserves the labial onset as a semi-vowel [v]. Liquid stability is noted in the Southern node. \\ 
\addlinespace[10pt]
\textbf{*KWAN} & svan & nay & khyi & \small PED *kwan. Vedic 'śvan' preserves the [K-W-N] cluster. Tibetan 'khyi' shows the velar radiation [kh]. Tamil 'nāy' (from *ñāy) shows the nasal shift of the onset, a common Dravidian palatalization. \\ 
\addlinespace[10pt]
\textbf{*SI-K} & yuka & pen & shig & \small PED *si-k. Tibetan 'shig' preserves the sibilant-velar [S-K] skeleton. Vedic 'yūka' shows the velar [k]. Tamil 'peṇ' is likely a distinct development, but the [S-K] alignment in the North/East is diagnostic. \\ 
\addlinespace[10pt]
\textbf{*MAR} & vrksa & maram & shing & \small PED *mar. Tamil 'maram' is the primary reflex. Tibetan 'shing' (from *si-ng) shows the nasalization of the coda. Vedic 'vṛkṣa' is IE, but 'mūla' (root) or local substrate 'mara' (in compounds) hints at the PED presence. \\ 
\addlinespace[10pt]
\textbf{*WI-T} & bija & vittu & sabon & \small PED *wi-t. Tamil 'vittu' preserves the [W-T] skeleton. Vedic 'bīja' shows the labial [b] and palatal [j < t]. Tibetan 'sabon' is a compound but retains the labial-nasal core. \\ 
\addlinespace[10pt]
\textbf{*LA-P} & parna & ilai & lo & \small PED *la-p. Tamil 'ilai' (from *ila-p?) preserves the liquid. Tibetan 'lo' is the reduced form. Vedic 'parṇa' shows the labial and liquid (p-r), with a nasal suffix. This root demonstrates 'Liquid Drift'. \\ 
\addlinespace[10pt]
\textbf{*RU-T} & mula & ver & rtsa & \small PED *ru-t. Tamil 'vēr' preserves the liquid. Tibetan 'rtsa' preserves the [R-T] skeleton with affrication. Vedic 'mūla' reflects a liquid shift (r > l). The dental coda follows the 'Dental-Palatal Invariant'. \\ 
\addlinespace[10pt]
\textbf{*PA-K} & tvac & tol & pags & \small PED *pa-k. Tibetan 'pags' and Tamil 'tōl' (from *p-tōl?) align on the labial/velar axis. Vedic 'tvac' shows the dental-velar cluster, reflecting the complexity of the skin/covering root in the proto-language. \\ 
\addlinespace[10pt]
\textbf{*SA-N} & mamsa & un & sha & \small PED *sa-n. Vedic 'māṃsa' preserves the sibilant and nasal. Tibetan 'sha' reflects the palatalized sibilant. Tamil 'ūṇ' (food/meat) preserves the nasal coda. Sibilant-to-palatal shift is a key new law. \\ 
\addlinespace[10pt]
\textbf{*KAR} & asrj & kuruti & khrag & \small PED *kar. Tamil 'kuruti' preserves the [K-R] skeleton. Tibetan 'khrag' shows the velar radiation [kh] and liquid [r]. Vedic 'asṛj' is IE, but the [K-R] root is ubiquitous in Dravidian and Tibeto-Burman. \\ 
\addlinespace[10pt]
\textbf{*RU-S} & asthi & elumpu & rus & \small PED *ru-s. Tibetan 'rus' is the primary reflex. Tamil 'elumpu' (from *el-u-mpu?) shows liquid/nasal clusters. Vedic 'asthi' shows the sibilant-dental [S-T] which often alternates with [R] in PED. \\ 
\addlinespace[10pt]
\textbf{*KA-M} & srnga & kompu & rwa & \small PED *ka-m. Tamil 'kōmpu' preserves the [K-M] skeleton. Vedic 'śṛṅga' shows the velar [g] and nasal [ṅ]. Tibetan 'rwa' is a distinct development, but the South/West alignment on the velar radiation is strong. \\ 
\addlinespace[10pt]
\textbf{*TA-L} & sirsa & talai & mgo & \small PED *ta-l. Tamil 'talai' is the perfect reflex of the 'Dental-Palatal Invariant' and 'Liquid Stability'. Vedic 'śīrṣa' shows the palatalized sibilant [ś < t]. Tibetan 'mgo' is a distinct velar-prefixed form. \\ 
\addlinespace[10pt]
\textbf{*KA-N} & karna & cevi & rna & \small PED *ka-n. Vedic 'karṇa' preserves the [K-N] skeleton. Tibetan 'rna' preserves the nasal. Tamil 'cevi' shows the palatalization of the velar initial (*k > c), confirming the 'Velar Radiation' law. \\ 
\addlinespace[10pt]
\textbf{*KA-N} & aksi & kan & mig & \small PED *ka-n. Tamil 'kaṇ' preserves the velar and nasal. Vedic 'akṣi' preserves the velar [k]. Tibetan 'mig' (from *m-ik) shows the velar coda and a nasal prefix, a classic 'Radial Nasal Shift' example. \\ 
\addlinespace[10pt]
\textbf{*NA-S} & nasa & mukku & sna & \small PED *na-s. Vedic 'nasa' and Tibetan 'sna' preserve the [N-S] skeleton. Tamil 'mūkku' shows the nasal shift [n > m] and velar radiation [k], linking 'nose' to the breathing 'velar' sound. \\ 
\addlinespace[10pt]
\textbf{*KA-W} & asya & vay & kha & \small PED *ka-w. Tibetan 'kha' preserves the velar radiation. Tamil 'vāy' preserves the labial/glide coda. Vedic 'āsya' is distinct, but the K-W alignment in the East/South is a diagnostic feature of the PED 'opening' root. \\ 
\addlinespace[10pt]
\textbf{*PA-L} & danta & pal & so & \small PED *pa-l. Tamil 'pal' is the primary reflex. Vedic 'danta' is IE but 'phala' (point) might be related. Tibetan 'so' reflects the sibilant shift. Liquid [l] stability is the key marker here. \\ 
\addlinespace[10pt]
\textbf{*LA-K} & jihva & nakku & lce & \small PED *la-k. Tamil 'nākku' shows the nasalized onset and velar coda. Tibetan 'lce' preserves the liquid [l]. Vedic 'jihvā' reflects the velar [h < k] and glide [v], aligning with the 'Velar Radiation' law. \\ 
\addlinespace[10pt]
\textbf{*KA-L} & pada & kal & rkang & \small PED *ka-l. Tamil 'kāl' preserves the [K-L] skeleton. Tibetan 'rkang' preserves the velar and adds a nasal coda. Vedic 'pada' is IE, but 'kula' (slope/base) provides a velar-liquid correlate. \\ 
\addlinespace[10pt]
\textbf{*PU-S} & janu & mulankal & pus & \small PED *pu-s. Tibetan 'pus' is the primary reflex. Tamil 'muḻaṅkāl' is a compound but contains the nasal-velar cluster. Vedic 'jānu' is IE, but the [P-S] root in Tibetan is a robust PED candidate. \\ 
\addlinespace[10pt]
\textbf{*KA-T} & hasta & kai & lag & \small PED *ka-t. Vedic 'hasta' preserves the [K-T] skeleton (h < k). Tamil 'kai' preserves the velar onset. Tibetan 'lag' shows the velar coda [g], following the 'Velar Radiation' across the limb terms. \\ 
\addlinespace[10pt]
\textbf{*WA-T} & udara & vayiru & grod & \small PED *wa-t. Tamil 'vayiṟu' preserves the [W-T/R] skeleton. Vedic 'udara' (from *wa-dar?) preserves the dental-liquid. Tibetan 'grod' shows the dental coda and velar onset radiation. \\ 
\addlinespace[10pt]
\textbf{*NI-NG} & hrdaya & neñcu & snying & \small PED *ni-ng. Tibetan 'snying' and Tamil 'neñcu' are perfect cognates following the 'Radial Nasal Shift' and 'Palatalization'. Vedic 'hṛdaya' shows the velar [h < k/n?] and dental [d]. \\ 
\addlinespace[10pt]
\textbf{*PA-K} & pa & kuti & thung & \small PED *pa-k. Vedic 'pā' preserves the labial onset. Tamil 'kuṭi' and Tibetan 'thung' show the velar and dental radiations respectively, suggesting a complex root for 'ingestion'. \\ 
\addlinespace[10pt]
\textbf{*A-N} & ad & un & za & \small PED *a-n. Tamil 'uṇ' preserves the nasal. Vedic 'ad' preserves the dental coda. Tibetan 'za' (from *ja < *n-ja?) reflects the palatalization of the nasal-dental cluster. \\ 
\addlinespace[10pt]
\textbf{*KA-N} & drs & kan & mthong & \small PED *ka-n. Tamil 'kaṇ' (eye/see) is the primary reflex. Tibetan 'mthong' shows the nasal-velar cluster. Vedic 'dṛś' is IE, but 'khyā' (to see/be known) preserves the velar radiation. \\ 
\addlinespace[10pt]
\textbf{*KA-L} & sru & kel & thos & \small PED *ka-l. Tamil 'kēḷ' preserves the [K-L] skeleton. Tibetan 'thos' shows the dental radiation [th]. Vedic 'śru' (from *klu) preserves the liquid and velar-derived sibilant. \\ 
\addlinespace[10pt]
\textbf{*SA-N} & jña & ari & shes & \small PED *sa-n. Tibetan 'shes' and Vedic 'jñā' (via *snyā?) reflect the sibilant/palatal onset. Tamil 'ari' preserves the liquid/vowel core. This root marks the cognitive boundary of the PED lexicon. \\ 
\addlinespace[10pt]
\textbf{*TA-L} & stha & nil & lang & \small PED *ta-l. Vedic 'sthā' preserves the dental onset [t]. Tamil 'nil' preserves the liquid coda [l] and nasal onset. Tibetan 'lang' shows the liquid and nasal radiation. \\ 
\addlinespace[10pt]
\textbf{*NYI} & surya & ñayiru & nyi & \small PED *nyi. Tibetan 'nyi' and Tamil 'ñāyiru' (from *ñāy-ir) are perfect reflexes of the palatal nasal. Vedic 'sūrya' preserves the glide and liquid, with the sibilant derived from the palatal onset. \\ 
\addlinespace[10pt]
\end{longtable}